\[ %%%%%%%%%%%%%%%%%%%%%%%%%%%%%%%%%%%%%%%%%%%%%%%%%%%%%%%%%%%%%%%%%%%%%%%%%%%%%%%% \subsection{Computational Efficiency} \label{sec:computational_efficiency} %%%%%%%%%%%%%%%%%%%%%%%%%%%%%%%%%%%%%%%%%%%%%%%%%%%%%%%%%%%%%%%%%%%%%%%%%%%%%%%% In addition to predictive accuracy and physical consistency, computational efficiency is an important requirement for practical scientific simulation and digital twin applications. This section evaluates the computational characteristics of PHYSGEN and compares them with representative baseline approaches. The objective is to determine whether the proposed physics-guided graph architecture provides improved performance without excessive computational overhead. %%%%%%%%%%%%%%%%%%%%%%%%%%%%%%%%%%%%%%%%%%%%%%%%%%%%%%%%%%%%%%%%%%%%%%%%%%%%%%%% \subsubsection{Evaluation Metrics} %%%%%%%%%%%%%%%%%%%%%%%%%%%%%%%%%%%%%%%%%%%%%%%%%%%%%%%%%%%%%%%%%%%%%%%%%%%%%%%% Computational efficiency is evaluated using four main metrics: \begin{enumerate} \item number of trainable parameters; \item training computational cost; \item inference time; \item memory consumption. \end{enumerate} The computational efficiency score is defined as: \begin{equation} C_{eff} = \frac{Performance} {Computational\ Cost}. \end{equation} %%%%%%%%%%%%%%%%%%%%%%%%%%%%%%%%%%%%%%%%%%%%%%%%%%%%%%%%%%%%%%%%%%%%%%%%%%%%%%%% \subsubsection{Model Complexity} %%%%%%%%%%%%%%%%%%%%%%%%%%%%%%%%%%%%%%%%%%%%%%%%%%%%%%%%%%%%%%%%%%%%%%%%%%%%%%%% The number of trainable parameters is calculated as: \begin{equation} N_p = \sum_i |\theta_i|. \end{equation} Although PHYSGEN contains multiple functional modules, the architecture is designed to share latent representations between physical constraints, graph evolution, and stability regulation. This avoids unnecessary duplication of computational components. %%%%%%%%%%%%%%%%%%%%%%%%%%%%%%%%%%%%%%%%%%%%%%%%%%%%%%%%%%%%%%%%%%%%%%%%%%%%%%%% \subsubsection{Inference Complexity} %%%%%%%%%%%%%%%%%%%%%%%%%%%%%%%%%%%%%%%%%%%%%%%%%%%%%%%%%%%%%%%%%%%%%%%%%%%%%%%% For graph-based physical systems, the computational complexity of message passing is approximately: \begin{equation} O(|E|d), \end{equation} where: \begin{itemize} \item $|E|$ is the number of graph edges; \item $d$ is the latent feature dimension. \end{itemize} The complexity of one evolution step is: \begin{equation} H_{t+1} = \Psi_\theta(H_t,E_t). \end{equation} Therefore, long-horizon prediction complexity scales as: \begin{equation} O(K|E|d), \end{equation} where $K$ represents the prediction horizon. %%%%%%%%%%%%%%%%%%%%%%%%%%%%%%%%%%%%%%%%%%%%%%%%%%%%%%%%%%%%%%%%%%%%%%%%%%%%%%%% \subsubsection{Scaling with Graph Size} %%%%%%%%%%%%%%%%%%%%%%%%%%%%%%%%%%%%%%%%%%%%%%%%%%%%%%%%%%%%%%%%%%%%%%%%%%%%%%%% To evaluate scalability, experiments are performed with increasing graph sizes: \begin{equation} N\in \{10^2,10^3,10^4\}. \end{equation} The evaluation measures: \begin{itemize} \item inference latency; \item GPU memory usage; \item prediction accuracy degradation. \end{itemize} The graph-based formulation allows PHYSGEN to scale with sparse physical interactions rather than dense state-space representations. %%%%%%%%%%%%%%%%%%%%%%%%%%%%%%%%%%%%%%%%%%%%%%%%%%%%%%%%%%%%%%%%%%%%%%%%%%%%%%%% \subsubsection{Comparison with Baseline Models} %%%%%%%%%%%%%%%%%%%%%%%%%%%%%%%%%%%%%%%%%%%%%%%%%%%%%%%%%%%%%%%%%%%%%%%%%%%%%%%% The expected computational comparison is summarized in Table~\ref{tab:computational_efficiency}. \begin{table}[h] \centering \caption{ Computational efficiency comparison. } \label{tab:computational_efficiency} \begin{tabular}{lcccc} \hline \textbf{Model} & \textbf{Parameters} & \textbf{Inference Cost} & \textbf{Memory} & \textbf{Scalability} \\ \hline MLP & Low & Low & Low & Low \\ LSTM & Medium & Medium & Medium & Medium \\ GNN & Medium & Medium & Medium & High \\ GNS & High & High & High & High \\ PINN & Medium & High & Medium & Medium \\ FNO & High & Medium & High & High \\ PHYSGEN & Medium & Medium & Medium & High \\ \hline \end{tabular} \end{table} %%%%%%%%%%%%%%%%%%%%%%%%%%%%%%%%%%%%%%%%%%%%%%%%%%%%%%%%%%%%%%%%%%%%%%%%%%%%%%%% \subsubsection{Accuracy-Efficiency Trade-off} %%%%%%%%%%%%%%%%%%%%%%%%%%%%%%%%%%%%%%%%%%%%%%%%%%%%%%%%%%%%%%%%%%%%%%%%%%%%%%%% To evaluate the balance between prediction quality and computational cost, we measure: \begin{equation} R_{AE} = \frac{ 1/RMSE } { T_{inference} }. \end{equation} A desirable model should achieve: \begin{equation} R_{AE}\rightarrow max. \end{equation} PHYSGEN aims to achieve improved accuracy-efficiency trade-off through: \begin{itemize} \item sparse graph representation; \item shared latent physical states; \item adaptive message passing; \item stability-aware rollout. \end{itemize} %%%%%%%%%%%%%%%%%%%%%%%%%%%%%%%%%%%%%%%%%%%%%%%%%%%%%%%%%%%%%%%%%%%%%%%%%%%%%%%% \subsubsection{Visualization} %%%%%%%%%%%%%%%%%%%%%%%%%%%%%%%%%%%%%%%%%%%%%%%%%%%%%%%%%%%%%%%%%%%%%%%%%%%%%%%% The computational scaling analysis is presented in: \begin{figure}[h] \centering \includegraphics[width=0.9\linewidth]{figures/computational_scaling.pdf} \caption{ Computational scaling of PHYSGEN compared with baseline models. The figure shows inference cost and memory requirements as a function of graph size. } \label{fig:computational_scaling} \end{figure} %%%%%%%%%%%%%%%%%%%%%%%%%%%%%%%%%%%%%%%%%%%%%%%%%%%%%%%%%%%%%%%%%%%%%%%%%%%%%%%% \subsubsection{Discussion} %%%%%%%%%%%%%%%%%%%%%%%%%%%%%%%%%%%%%%%%%%%%%%%%%%%%%%%%%%%%%%%%%%%%%%%%%%%%%%%% The computational analysis demonstrates that physics-guided learning does not necessarily require prohibitive computational complexity. By representing physical systems as adaptive graphs, PHYSGEN avoids the limitations of dense discretization approaches and provides efficient scaling for large dynamical systems. This property is particularly important for real-time digital twins, where prediction speed and reliability are equally important. %%%%%%%%%%%%%%%%%%%%%%%%%%%%%%%%%%%%%%%%%%%%%%%%%%%%%%%%%%%%%%%%%%%%%%%%%%%%%%%% \subsubsection{Summary} %%%%%%%%%%%%%%%%%%%%%%%%%%%%%%%%%%%%%%%%%%%%%%%%%%%%%%%%%%%%%%%%%%%%%%%%%%%%%%%% The computational efficiency evaluation shows that PHYSGEN provides a balanced combination of predictive performance, physical reliability, and scalability. The results indicate that the proposed framework is suitable for large-scale scientific simulation and future digital twin applications. \] 
